\bookchapter{الفصل 15: علم البطارية وتحسين مدة الاستخدام}

\chapterintro{البطارية هي “قلب الطاقة” لهاتفك. يأخذك هذا الفصل لفهم طباع بطاريتك الصغيرة، وتعلّم تقنيات الشحن وإطالة العمر بشكل علمي، لتودع قلق نفاد الطاقة إلى الأبد.}

\booksection{أولاً: كيف تعمل البطارية؟ دورات الشحن ومستوى الصحة}

بطارية الهاتف (ليثيوم-أيون أو ليثيوم-بوليمر) تشبه إسفنجة الطاقة: الشحن هو امتصاص الماء (تخزين الكهرباء)، والتفريغ هو عصر الماء (استهلاك الكهرباء). هناك خصائص تستحق التذكر:

\begin{itemize}
  \item \textbf{الشحن عند الحاجة}: على عكس البطاريات القديمة، ليس لها تأثير ذاكرة، ولا داعي لتفريغها بالكامل قبل الشحن.
  \item \textbf{تخاف من الحرارة}: الحرارة المرتفعة هي القاتل الأول للبطارية.
  \item \textbf{تتقدّم في العمر}: مع الاستخدام المطوّل، تنخفض السعة القصوى، وهذا أمر طبيعي.
\end{itemize}

\booksubsection{المؤشرات الجوهرية: عدد الدورات ومستوى الصحة}

\begin{itemize}
  \item الدورة الواحدة لا تعني شحنة واحدة. مثلاً: إذا استخدمت اليوم 60\% ثم شحنته بالكامل، وغداً استخدمت 40\% ثم شحنته بالكامل، فإن مجموع 100\% من الاستهلاك يُحتسب دورة واحدة.
  \item مستوى صحة البطارية هو نسبة السعة القصوى الحالية مقارنة بالحالة الجديدة. عندما ينخفض عن 80\% ستشعر بوضوح أن مدة الاستخدام قد قصرت.
  \item \textbf{كيف تتحقق؟}iPhone: الإعدادات > البطارية > صحة البطارية والشحن؛ أندرويد: ابحث عن "صحة البطارية" في الإعدادات.
\end{itemize}

\booksection{ثانياً: أساطير الشحن؟ الشحن السريع واللاسلكي والقاعدة الذهبية}

\begin{itemize}
  \item \textbf{هل الشحن السريع يضر البطارية؟}باستخدام ملحقات الشحن السريع الأصلية أو المعتمدة، يكون الضرر ضئيلاً. الهاتف ورأس الشحن مزوّدان بتحكم في الحرارة، ويخفّضان الطاقة تلقائياً عند ارتفاع الحرارة.
  \item \textbf{الشحن اللاسلكي}: مريح لكنه أقل كفاءة ويولّد حرارة أكبر. مناسب للشحن طوال الليل.
  \item \textbf{القاعدة الذهبية للحفاظ على الصحة}: حافظ على مستوى البطارية بين 20\%-80\% قدر الإمكان، ولا تدعها تجوع (0\%) أو تشبع (100\%) باستمرار.
\end{itemize}

\begin{tcbinfo}{وظيفة "حد شحن 80\%" في iPhone — تحويل القاعدة الذهبية إلى إعداد نظام}
سلسلة iPhone 15/16 (iOS 17 وأحدث) حوّلت “قاعدة 20\%-80\%” إلى وظيفة نظام. عند التفعيل، يتوقف الهاتف عن الشحن تلقائياً عند 80\%، مما يبطئ شيخوخة البطارية بشكل كبير. مسار الإعداد:

\texttt{الإعدادات > البطارية > صحة البطارية والشحن > تحسين الشحن > حد 80\%}

معاني الخيارات الثلاثة: (1)\textbf{غير محدود}— يشحن إلى 100\%؛ (2)\textbf{تحسين شحن البطارية}— يتعلم روتينك، ويؤخر الشحن إلى 100\% فقط في الفترة الأخيرة قبل الفجر؛ (3)\textbf{حد 80\%}— يتوقف مباشرة عند 80\%، الأكثر حماية للبطارية. لمن يتوفر لديه شحن في المنزل/المكتب بشكل دائم، يُنصح بشدة باختيار حد 80\%؛ لمن يسافر كثيراً ويحتاج بطارية ممتلئة عند الخروج، يمكن اختيار تحسين شحن البطارية.

\end{tcbinfo}

\begin{tcbinfo}{شاحن GaN — الرفيق المثالي لمن يسافر}
ربما لاحظت أن رؤوس الشحن السريع الحالية أصبحت أصغر فأصغر، والبطل وراء ذلك هو\textbf{GaN (نتريد الغاليوم)}هذه المادة أشباه الموصلات الجديدة. مقارنة بالشواحن التقليدية القائمة على السيليكون، لشواحن GaN مزايا واضحة:

\begin{itemize}
  \item \textbf{حجم صغير}: بنفس القدرة، يمكن لشاحن GaN أن يكون نصف أو ثلث حجم الشاحن التقليدي، ورأس GaN بقدرة 65W أصغر من رأس Apple القديم 5W.
  \item \textbf{كفاءة عالية وحرارة منخفضة}: خسارة التبديل في GaN أقل بكثير من السيليكون، وكفاءة تحويل الطاقة تصل إلى أكثر من 95\%، ولا يسخن كثيراً أثناء الشحن.
  \item \textbf{مناسب للسفر}: شاحن GaN متعدد المنافذ بقدرة 65W أو 100W واحد، يمكنه شحن الهاتف واللوحي والحاسوب في آن واحد، أخف بكثير من حمل ثلاث "قرميدات".
  \item \textbf{نصيحة الشراء}: ابحث عن شواحن GaN الداعمة لبروتوكول PD (65W فأكثر يمكنه شحن الحاسوب)، والعلامات المعروفة مثل Baseus وAnker وNubia وXiaomi وغيرها جميعها جيدة. انتبه للتأكد من تطابق كابل البيانات (القدرة العالية تتطلب كابل Type-C برقاقة E-Marker).
\end{itemize}

\end{tcbinfo}

\booksection{ثالثاً: تحليل شامل لبروتوكولات الشحن السريع: لماذا لا يشحن شاحنك بسرعة؟}

ربما واجهت هذا الموقف: هاتفك يدعم الشحن السريع، لكن عند تغيير الشاحن أصبح بطيئاً جداً. ذلك لأن الشحن السريع ليس مجرد “طاقة عالية تكفي”، بل يتطلب أن “يتحدث رأس الشحن وكابل البيانات والهاتف اللغة نفسها” — وهذه هي بروتوكولات الشحن السريع.

\booksubsection{المبدأ الأساسي للشحن السريع}

قدرة الشحن = الجهد × التيار (P = U × I). لزيادة القدرة، إمّا رفع الجهد أو رفع التيار أو كليهما. وقد اختارت البروتوكولات المختلفة مسارات تقنية مختلفة.

سير عمل الشحن السريع:

\begin{enumerate}
  \item \textbf{مرحلة المصافحة}: يتواصل الهاتف ورأس الشحن عبر كابل البيانات، ويتبادلان معلومات البروتوكولات والقدرات المدعومة.
  \item \textbf{مرحلة التفاوض}: يحدّد الطرفان مستوى القدرة القصوى المشترك المدعوم.
  \item \textbf{مرحلة الشحن}: يخرج رأس الشحن الجهد والتيار المتفق عليهما، ويقوم الهاتف بإدارة الطاقة.
  \item \textbf{مرحلة التعديل}: مع زيادة الشحن، يخفّض الهاتف تدريجياً قدرة الشحن لحماية البطارية.
\end{enumerate}

\booksubsection{تفاصيل بروتوكولات الشحن السريع الرئيسية}

\textbf{USB PD (Power Delivery)}— ملك البروتوكولات العامة

\begin{itemize}
  \item \textbf{الميزات}: معيار رسمي من منظمة USB-IF، الأوسع انتشاراً. يدعم PD 3.0 حتى 100 واط (20V/5A)، ويمتد PD 3.1 إلى 240 واط (48V/5A).
  \item \textbf{مستويات الجهد}: 5V، 9V، 15V، 20V، ويدعم PD 3.1 أيضاً 28V، 36V، 48V.
  \item \textbf{التوافق}: جميع أجهزة iPhone، وGoogle Pixel، وبعض أجهزة أندرويد الرائدة، والحاسبات المحمولة، وجهاز ألعاب Switch تدعم بروتوكول PD.
  \item \textbf{المزايا}: توافق عالٍ، رأس شحن واحد لشحن أجهزة متعددة؛ يدعم PPS (مصدر قابل للبرمجة) لضبط دقيق للجهد والتيار.
  \item \textbf{العيوب}: يتطلب واجهة Type-C وشريحة مخصصة، التكلفة أعلى.
\end{itemize}

\textbf{Qualcomm QC (Quick Charge)}— رائد معسكر أندرويد

\begin{itemize}
  \item \textbf{QC 3.0}: يدعم جهداً قابلاً للضبط من 3.6V-20V بزيادات 0.2V، حتى 18 واط.
  \item \textbf{QC 4.0/4+}: متوافق مع بروتوكول USB PD، يدعم حتى 27 واط، ويضيف تقنية الشحن ثنائي المسار.
  \item \textbf{QC 5.0}: يدعم أكثر من 100 واط، 0-50\% في 5 دقائق فقط، مع تحسين كبير في إدارة الحرارة.
  \item \textbf{التوافق}: معظم الهواتف التي تستخدم معالجات Snapdragon من Qualcomm تدعمه، لكن العلامات التجارية عادةً تعطي الأولوية لبروتوكولاتها الخاصة.
\end{itemize}

\textbf{OPPO VOOC / OnePlus WARP / Realme Dart}— ممثلو الجهد المنخفض والتيار العالي

\begin{itemize}
  \item \textbf{المسار التقني}: جهد منخفض وتيار عالٍ (يبدأ من 5V/4A)، حرارة أقل، كفاءة أعلى.
  \item \textbf{VOOC 1.0}: 20 واط (5V/4A)، أُطلق عام 2014، فتح عصر الشحن السريع.
  \item \textbf{SuperVOOC}: 50-65 واط، يعتمد تصميم خليتين متتاليتين.
  \item \textbf{SuperVOOC 2.0}: 65 واط، يشحن بطارية 4000mAh في 30 دقيقة.
  \item \textbf{الشحن الفائق 150 واط/240 واط}: أسرع تقنية شحن سريع في الإنتاج التجاري حالياً.
  \item \textbf{الميزات}: يتطلب رأس شحن وكابل بيانات مخصصين (بنقاط تلامس إضافية)، غير متوافق مع بروتوكولات أخرى.
\end{itemize}

\textbf{Huawei SuperCharge}— تطور الشحن السريع عالي الجهد

\begin{itemize}
  \item \textbf{FCP (Fast Charge Protocol)}: 9V/2A، 18 واط، متوافق مع QC 2.0.
  \item \textbf{SCP (Super Charge Protocol)}: يدعم قدرات أعلى، من 22.5 واط إلى 66 واط ثم 100 واط.
  \item \textbf{الخصائص التقنية}: يعتمد تقنية مضخة الشحن، يحوّل الجهد العالي إلى جهد منخفض وتيار عالٍ يدخل البطارية، بكفاءة تصل إلى 98\%.
  \item \textbf{التوافق}: مخصص لهواتف Huawei/Honor، وبعض الشواحن تدعم بروتوكول PD.
\end{itemize}

\textbf{شحن Xiaomi Surge السريع}— صعود البروتوكول الذاتي

\begin{itemize}
  \item \textbf{شريحة Surge P1}: شريحة إدارة الشحن السريع ذاتية التطوير من Xiaomi، تدعم الشحن السريع السلكي 120 واط.
  \item \textbf{شريحة Surge G1}: شريحة إدارة البطارية، تعمل مع P1 لتحقيق شحن سريع أكثر كفاءة.
  \item \textbf{الخصائص التقنية}: يعتمد مخطط خليتين متتاليتين + مضخة شحن، يدعم توافق بروتوكولات متعددة.
\end{itemize}

\textbf{شحن Apple السريع}— محافظ لكن مستقر

\begin{itemize}
  \item \textbf{البروتوكول}: يدعم فقط بروتوكول USB PD.
  \item \textbf{القدرة}: iPhone 8 وما بعده يدعم الشحن السريع، حتى حوالي 27 واط (يمكن لـ iPhone 15 Pro Max بلوغ 27+ واط).
  \item \textbf{الميزات}: القدرة محافظة نسبياً، لكن منحنى الشحن مُحسّن جيداً، والتحكم بالحرارة ممتاز.
  \item \textbf{متطلبات الملحقات}: يحتاج كابل USB-C to Lightning (أو USB-C to USB-C) ورأس شحن PD.
\end{itemize}

\booksubsection{جدول مقارنة بروتوكولات الشحن السريع}

\begin{center}
\begin{tabularx}{\linewidth}{|X|X|X|X|}
\hline
\textbf{البروتوكول} & \textbf{القدرة القصوى} & \textbf{المسار التقني} & \textbf{التوافق} \tabularnewline
\hline
USB PD 3.1 & 240 واط & جهد عالٍ/مستويات متعددة & الأوسع \tabularnewline
QC 5.0 & +100 واط & جهد عالٍ قابل للضبط & أجهزة Qualcomm \tabularnewline
VOOC/Dart & 240 واط & جهد منخفض وتيار عالٍ & مخصص لسلسلة OPPO \tabularnewline
SuperCharge & 100 واط & مضخة شحن & Huawei/Honor \tabularnewline
شحن Apple السريع & \textasciitilde{}27 واط & بروتوكول PD & iPhone/iPad \tabularnewline
\hline
\end{tabularx}
\end{center}

\booksubsection{بروتوكولات الشحن اللاسلكي}

للشحن اللاسلكي معاييره وبروتوكولاته الخاصة:

\begin{itemize}
  \item \textbf{معيار Qi}: وضعته WPC (تحالف الشحن اللاسلكي)، الأكثر شيوعاً. يدعم Qi الأساسي من 5W إلى 15W.
  \item \textbf{معيار Qi2}: صدر عام 2023، يُدخل المحاذاة المغناطيسية (مشابه لـ MagSafe من Apple)، يدعم 15W+.
  \item \textbf{البروتوكولات الخاصة}: لكل علامة تجارية شحنها اللاسلكي عالي السرعة، مثل الشحن اللاسلكي 50 واط من Xiaomi و50 واط من Huawei، ويتطلب شواحن مخصصة.
\end{itemize}

\booksubsection{دليل شراء الشحن السريع}

\begin{itemize}
  \item \textbf{تحقق من توافق البروتوكول}: عند شراء شاحن، تأكد من دعمه لبروتوكول الشحن السريع لهاتفك. عادةً ما تذكر الشواحن من الأطراف الثالثة البروتوكولات المدعومة.
  \item \textbf{تحقق من كابل البيانات}: الشحن السريع عالي القدرة يتطلب كابل بيانات مخصصاً. الكابلات العادية قد لا تدعم أكثر من 60 واط.
  \item \textbf{تحقق من الشهادات}: يُنصح مستخدمو Apple باختيار كابلات معتمدة MFi؛ ويختار مستخدمو أندرويد الملحقات الأصلية أو من أطراف ثالثة معروفة.
  \item \textbf{انتبه للحرارة}: ارتفاع حرارة الهاتف أثناء الشحن السريع أمر طبيعي، لكن إذا أصبح ساخناً جداً فهذا غير طبيعي، يُنصح بإيقاف الشحن أو تغيير الشاحن.
\end{itemize}

\begin{tcbwarn}{تنبيهات سلامة الشحن السريع}
\begin{itemize}
  \item تجنب لعب الألعاب الكبيرة أثناء الشحن، فهو يزيد الحرارة.
  \item لا تستخدم شواحن وكابلات رديئة، فقد تؤدي إلى مخاطر سلامة.
  \item أزل غلاف الهاتف السميك أثناء الشحن لمساعدة التبريد.
  \item إذا ارتفعت حرارة الشاحن أو الهاتف بشكل غير طبيعي أو ظهرت رائحة غريبة، توقف عن الاستخدام فوراً.
\end{itemize}

\end{tcbwarn}

\booksection{أربعة: تكنولوجيا البطارية الجديدة: بطارية الأنود الكربوني السيليكوني}

إذا تابعت إطلاقات الهواتف الرائدة 2024-2026، أصبحت كلمة “الأنود الكربوني السيليكوني” مرادفاً للبطاريات المتقدمة. هذه ترقية مهمة لمادة أنود بطارية الليثيوم، والمبدأ ليس معقداً:

\begin{itemize}
  \item \textbf{الأنود الجرافيت التقليدي}: بطاريات الهواتف تستخدم الجرافيت كأنود دائماً، التقنية ناضجة ومستقرة، لكن كثافة الطاقة وصلت قريباً من الحد النظري (حوالي 372mAh/g).
  \item \textbf{إضافة مادة السيليكون}: السعة النوعية النظرية للسيليكون حوالي 10 أضعاف الجرافيت (حوالي 4200mAh/g). إضافة كمية قليلة من السيليكون إلى الجرافيت لصنع "أنود مركب كربوني سيليكوني"، يمكن أن يرفع كثافة طاقة البطارية بأكملها بشكل ملحوظ.
  \item \textbf{تحسين كثافة الطاقة}: بنفس الحجم، يمكن زيادة سعة بطارية الأنود الكربوني السيليكوني بحوالي\textbf{15\%-20\%}. وهذا هو السبب الرئيسي لقدرة الهواتف الرائدة 2024-2026 على حزم بطاريات كبيرة 5500-6000mAh في هيكل بسمك 7-9mm.
\end{itemize}

\begin{tcbinfo}{أي الموديلات استخدمتها؟}
\begin{itemize}
  \item \textbf{Honor}: سلسلة Magic رائدة في الاستخدام التجاري للأنود الكربوني السيليكوني، Magic6 Pro وسلسلة Magic7 وMagic V3 القابل للطي جميعها تستخدمه.
  \item \textbf{Xiaomi}: Xiaomi 14 Ultra وسلسلة Xiaomi 15 تحمل بطارية أنود كربوني سيليكوني من الجيل الجديد.
  \item \textbf{vivo}: سلسلة X100 Ultra وX200 Pro تستخدم حل الأنود الكربوني السيليكوني عالي الكثافة.
  \item \textbf{Huawei/OPPO وغيرها}: الرائدة 2025-2026 تتبنى بشكل عام، والأنود الكربوني السيليكوني يتحول من "تقنية سوداء" إلى معيار أساسي للهواتف الرائدة.
\end{itemize}

\end{tcbinfo}

\textbf{التأثير العملي على المستخدم}: أكبر فائدة للأنود الكربوني السيليكوني هي “سعة بطارية أكبر بنفس الحجم” — يمكن جعل الهاتف أرق مع الحفاظ على البطارية؛ أو بقاء سماكة الهاتف كما هي لكن زيادة البطارية من 4500mAh إلى 5500mAh. بعبارة أخرى، لم تعد مضطراً لدفع ثمن “السماكة والوزن” مقابل البطارية. تجدر الإشارة إلى أن بطاريات الأنود الكربوني السيليكوني لا تزال بحاجة للتعاون مع رقائق الشحن السريع في إدارة درجات الحرارة المنخفضة والشحن الكامل، وهذا أيضاً أحد أسباب توصية الشركات بالحفاظ على شحن 80\% على المدى الطويل.

\booksection{خمسة: هذه هي نصائح توفير الطاقة الحقيقية}

\begin{itemize}
  \item \textbf{الشاشة (ملك استهلاك الطاقة)}: اخفض السطوع، قصّر وقت القفل التلقائي. هواتف OLED توفر الكثير من الطاقة في الوضع الداكن.
  \item \textbf{إشارة ضعيفة}: في المصاعد والأقبية، يبحث الهاتف بجدية عن إشارة، مما يستهلك الطاقة بسرعة.
  \item \textbf{5G}: إذا كانت الإشارة ضعيفة، بدّل يدوياً إلى 4G، وستتحسن مدة الاستخدام كثيراً.
  \item \textbf{خدمات الموقع}: للتطبيقات غير الملاحية، اضبط صلاحية الموقع إلى "السماح فقط أثناء الاستخدام".
  \item \textbf{تقييد النشاط في الخلفية}: iPhone في الإعدادات > عام > تحديث التطبيقات في الخلفية؛ أندرويد في إعدادات بطارية التطبيقات.
\end{itemize}

\begin{tcbpractice}{لنتحقق: أجرِ فحصاً صحياً لبطاريتك}
\begin{enumerate}
  \item تحقق من مستوى الصحة: ابحث عن وسجّل السعة القصوى لبطاريتك.
  \item حلّل تفاصيل الاستهلاك: انظر إلى الرسوم البيانية وقائمة التطبيقات في إعدادات البطارية، واكتشف المذنب الرئيسي للاستهلاك.
  \item ضع خطة: قيّد التطبيقات التي تستهلك بشكل غير طبيعي في الخلفية.
\end{enumerate}

\end{tcbpractice}
